\documentclass{article}
\usepackage{/globals/math}
\begin{document}
\section{Brüche}
Ein Bruch ist eine Division eines \emph{Zählers} durch einen \emph{Nenner}.
\[
 a : b = \frac{a}{b}
\]
Jede Zahl kann als Bruch dargestellt werden.
\[
 a = \frac{a}{1}
\]

\subsection{Erweitern}
Ein Bruch kann \emph{erweitern} werden; sowohl der Zähler als auch der Nenner wird mit der gleichen Zahl multipliziert. Dabei verändert sich der Wert des Bruchs nicht. Ist hier $n<0$ wird es auch kürzen genannt.
\[
 \frac{a}{b} = \frac{a*n}{b*n} \dtext{wenn} n \ne 0
\]
Soll ein Bruch in einer Aufgabe gekürzt werden, so soll er so weit gekürzt werden, dass er nicht weiter gekürzt werden kann, weil des größte gemeinsame Teiler eins ist.

\subsection{Addition und Subtraktion}
Zwei Brüche können addiert oder subtrahiert werden wenn sie den gleichen Nenner haben.
\[
 \frac{a}{c} + \frac{b}{c} = \frac{a+b}{c}
\]
Ist dies nicht der Fall, liegen zwei Brüche mit unterschiedlichen Nennern vor, so können sie erweitert werden dass sie einen gemeinsamen Nenner haben. Ist kein potentieler gemeinsamer Nenner offensichtlich, so kann jeder Bruch mit dem Nenner des anderen erweitert werden.
\[
 \frac{a}{b} + \frac{c}{d} = \frac{a*d}{b*d} + \frac{c*b}{d*b} = \frac{ad+bc}{bd}
\]

\subsection{Multiplikation und Division}
Bei der Multiplikation können einfach die Zähler miteinander und die Nenner miteinander multipliziert werden.
\[
 \frac{a}{b} * \frac{c}{d} = \frac{ac}{bd}
\]
Eine Division ist eine Multiplikation mit dem Kehrwert.
\[
 \frac{\frac{a}{b}}{\frac{c}{d}} = \frac{a}{b} * \frac{d}{c}
\]
Wird sowohl der Zähler und der Nenner eines Bruchs selbst als ein Bruch dargestellt oder das Kommutativgesetz genutzt, so kann ein Bruch als eine Multiplikation von zwei Brüchen dargestellt werden, mit einmal einem Zähler von eins und einmal einem Nenner von eins.
\[
 \frac{\frac{a}{1}}{\frac{b}{1}} = \frac{a}{1} * \frac{1}{b}
 \dtext{oder}
 \frac{a}{b} = \frac{1*a}{b} = a * \frac{1}{b} = \frac{a}{1} * \frac{1}{b}
\]
\end{document}